\documentclass{article}
\usepackage[margin=28mm]{geometry}
\usepackage{amsmath,amssymb}
\usepackage{xurl}
\usepackage[hidelinks]{hyperref}
% Shared commands for notes in docs/paper-gaps/.

\DeclareUnicodeCharacter{2080}{\ifmmode{}_{0}\else\textsubscript{0}\fi}
\DeclareUnicodeCharacter{2081}{\ifmmode{}_{1}\else\textsubscript{1}\fi}
\DeclareUnicodeCharacter{2082}{\ifmmode{}_{2}\else\textsubscript{2}\fi}
\DeclareUnicodeCharacter{2099}{\ifmmode{}_{n}\else\textsubscript{n}\fi}

\newcommand{\Lean}{\textsc{Lean}}
\ExplSyntaxOn
\NewDocumentCommand{\leanid}{m}{
  \ifmmode
    \textsf{\detokenize{#1}}
  \else
    \group_begin:
    \urlstyle{sf}
    \tl_set:Nn \l_tmpa_tl {#1}
    \tl_replace_all:Nnn \l_tmpa_tl {\_} {_}
    \exp_args:NV \path \l_tmpa_tl
    \group_end:
  \fi
}
\ExplSyntaxOff
\providecommand{\tr}{\operatorname{tr}}
\newcommand{\tnleanIssueBase}{https://github.com/LionSR/TNLean/issues/}
\newcommand{\tnleanPrBase}{https://github.com/LionSR/TNLean/pull/}
\newcommand{\tnleanDocBase}{https://sirui-lu.com/TNLean/docs/}
\newcommand{\ghissue}[1]{\href{\tnleanIssueBase#1}{GitHub issue \##1}}
\newcommand{\ghpr}[1]{\href{\tnleanPrBase#1}{PR \##1}}
\newcommand{\doclink}[1]{\href{\tnleanDocBase#1.html}{\path{#1}}}

% Dirac notation and the identity operator, as in blueprint/src/macros/common.tex.
% \providecommand keeps note-local definitions authoritative.
\usepackage{dsfont}
\providecommand{\ket}[1]{|#1\rangle}
\providecommand{\bra}[1]{\langle #1|}
\providecommand{\braket}[2]{\langle #1 | #2 \rangle}
\providecommand{\ketbra}[2]{|#1\rangle\!\langle #2|}
\providecommand{\Id}{\mathds{1}}
\providecommand{\one}{\mathds{1}}
\providecommand{\C}{\mathbb{C}}
\providecommand{\MN}[1]{M_{#1}(\C)}
\providecommand{\MD}{M_D(\C)}

% Verdict marker: \gapnote{<kind>}{<status>}. Kind is the policy
% classification (clarification, local-correction, scope-restriction,
% unfaithful, false-source, open-gap); status is open, wip (elimination
% underway), resolved, or historical. Machine-read by the site index;
% prints a verdict line.
\newcommand{\gapnote}[2]{\par\noindent\textbf{Verdict:} #1 (#2).\par}

\providecommand{\gapnote}[2]{\par\noindent\textbf{Verdict:} #1 (#2).\par}
\title{The Right Side of Hastings' Single-Cluster Bound}
\author{}
\date{}
\begin{document}
\maketitle
\gapnote{local-correction}{resolved}

\section{The printed step}
Hastings, \emph{Solving gapped Hamiltonians locally}~\cite{Hastings2006Solving},
\href{https://arxiv.org/abs/cond-mat/0508554v3}{arXiv:cond-mat/0508554v3},
source \path{proj3new.tex}, lines 1174--1184, defines $\rho(\beta,C)$ as the
sum of the terms of the power series of $e^{-\beta H}$ that contain every
$H_i$ with $i\in C$ and no $H_j$ with $j\in B(C)$, and states that for any
operator $O$
\begin{align}
 \operatorname{Tr}\bigl(O\rho(\beta,C)\bigr)
 \le\bigl[e^{J\beta}-1\bigr]^{n_C}e^{\beta\|H_C+H_{B(C)}\|}
 \operatorname{Tr}\bigl(Oe^{-\beta H}\bigr).
 \label{eq:hastings05_cluster_part_printed}
\end{align}

\section{The printed form fails}
Write $c\ge0$ for the constant on the right. For a Hermitian $O$,
applying~(\ref{eq:hastings05_cluster_part_printed}) to $O$ and to $-O$ forces
$\operatorname{Tr}(O\rho(\beta,C))=c\operatorname{Tr}(Oe^{-\beta H})$, so
$\rho(\beta,C)$ would be a multiple of $e^{-\beta H}$. A single term already
violates this. Take $H=h=\operatorname{diag}(1,0)$, $C=\{1\}$, $B(C)=\emptyset$
and $\beta>0$. Then $\rho(\beta,C)=e^{-\beta h}-1=\operatorname{diag}(e^{-\beta}-1,0)$,
and for $O=-\operatorname{diag}(1,0)$ the left side
of~(\ref{eq:hastings05_cluster_part_printed}) is $1-e^{-\beta}>0$ while the
right side is $-ce^{-\beta}\le0$.

\section{The correction}
The next display, equation (percb2) at lines 1186--1195, bounds
$Z^{-1}\operatorname{Tr}(O\rho(\beta,C))$ by $\|O\|$ times the same constant,
with $Z=\operatorname{Tr}e^{-\beta H}$. The intended intermediate bound is
therefore
\begin{align}
 \bigl|\operatorname{Tr}\bigl(O\rho(\beta,C)\bigr)\bigr|
 \le\bigl[e^{J\beta}-1\bigr]^{n_C}e^{\beta\|H_C+H_{B(C)}\|}\,\|O\|Z,
 \label{eq:hastings05_cluster_part_corrected}
\end{align}
which is what the formalization proves, with
$\|H_C+H_{B(C)}\|\le(n_C+n_{B(C)})J$. Equation (percb2) and everything
downstream are unaffected.
\bibliographystyle{unsrt}
\bibliography{references}
\end{document}
